Từ mục tiêu và phạm vi đã nêu, đồ án có năm đóng góp chính, tương ứng với các tầng của quy trình CausClass và cách đánh giá được triển khai trong các chương sau.

Thứ nhất, đồ án đề xuất CausClass, một khung thống nhất nối liền nhận diện hành vi lớp học, khám phá liên kết thời gian, tinh chỉnh đồ thị có hướng dẫn bởi tri thức và sinh báo cáo quan sát. Nhờ đó, video lớp học không chỉ được xử lý như tập hợp khung hình rời rạc, mà được chuyển thành chuỗi bằng chứng có cấu trúc: từ hộp phát hiện, quỹ đạo theo dõi, sự kiện vi mô, chuỗi thời gian, đồ thị đến báo cáo. Mỗi tầng đều có sản phẩm trung gian để người dùng kiểm tra lại.

Thứ hai, đồ án xây dựng và đánh giá mô-đun nhận diện hành vi tăng cường cơ chế chú ý dựa trên YOLO26n-MHSA. Biến thể này giữ nguyên phần lớn kiến trúc YOLO26n nhưng chèn một khối MHSA vào đường tổng hợp đặc trưng của phần đầu phát hiện để bổ sung ngữ cảnh không gian tầm xa. Thiết kế cục bộ này giúp so sánh với YOLO26n gốc rõ ràng hơn: khác biệt hiệu năng có thể quy chủ yếu cho mô-đun MHSA thay vì một tái thiết kế toàn bộ bộ phát hiện.

Thứ ba, đồ án giới thiệu chiến lược khám phá đồ thị có hướng dẫn bởi tri thức, trong đó LLM sinh các giả thuyết chỉnh sửa cấu trúc còn dữ liệu thời gian và AERCA quyết định việc giữ hay loại ứng viên. Cách thiết kế này kết hợp được tri thức ngữ nghĩa về lớp học với kiểm chứng định lượng, đồng thời tránh việc chấp nhận trực tiếp một đồ thị chỉ vì nó hợp lý về mặt ngôn ngữ.

Thứ tư, đồ án xây dựng một bộ kiểm chuẩn có thể tái lập gồm 100 đồ thị thời gian lớp học tổng hợp. Mỗi bộ kiểm thử có sáu biến hành vi, đồ thị chuẩn và chuỗi VAR(1) được sinh từ cấu trúc đã biết. Bộ kiểm chuẩn này cho phép đo Precision, Recall và F1 của các chiến lược tìm kiếm đồ thị trong điều kiện có đồ thị chuẩn, điều khó đạt được với video lớp học thật.

Thứ năm, đồ án kiểm chứng tính hiệu quả và khả năng diễn giải của CausClass thông qua thí nghiệm trên dữ liệu phát hiện hành vi, bộ kiểm chuẩn đồ thị thời gian tổng hợp và nghiên cứu tình huống trên video lớp học thật. Nghiên cứu tình huống cho thấy hệ thống có thể chuyển từ quan sát hành vi thô sang thống kê thời gian, đồ thị liên kết và báo cáo theo nguyên tắc nêu bằng chứng trước. Báo cáo cuối cùng được viết theo hướng thận trọng: diễn giải liên kết thời gian như quan hệ dự báo và yêu cầu người dùng kiểm tra lại video trước khi đưa ra nhận định sư phạm.

Các đóng góp trên phản ánh một quan điểm thiết kế xuyên suốt: hệ thống tự động phải có cơ chế tự giới hạn. Bộ phát hiện không được diễn giải thay cho giáo viên; đồ thị không được xem là bằng chứng nhân quả chắc chắn; và LLM không được quyết định cấu trúc chỉ dựa trên văn bản. Vì vậy, mỗi đóng góp đều đi kèm một cơ chế kiểm soát: đánh giá bộ phát hiện trên nhãn chung, bộ kiểm chứng định lượng cho đồ thị, bộ kiểm chuẩn có đồ thị chuẩn cho Stage 2 và báo cáo có ràng buộc bằng chứng cho phần diễn giải.